\documentclass{article}

\usepackage[preprint]{neurips_2026}

\usepackage[utf8]{inputenc}
\usepackage[T1]{fontenc}
\usepackage{hyperref}
\usepackage{url}
\usepackage{booktabs}
\usepackage{amsfonts}
\usepackage{amsmath}
\usepackage{nicefrac}
\usepackage{microtype}
\usepackage{xcolor}

\title{ Warm-Start Diffusion for Low-Resolution Probabilistic ERA5 Weather Forecasting}

\author{%
  Kabir Vats (Team 43)
}

\begin{document}

\maketitle

\section{Project Proposal: Track 1}

\subsection{Problem Background \& Motivation}

In this project, I will work on probabilistic forecasting on a low-resolution ERA5 dataset, using a $5^\circ$ global grid. The goal is to predict future weather states at short-to-medium lead times from recent atmospheric states. The low resolution is useful because it lets us test ideas quickly while still keeping the important structure of the problem: global spatial coupling, autoregressive rollout, uncertainty growth, and the need for physically reasonable samples.

Recent machine learning weather models have made this setting much more interesting. GraphCast showed that learned neural weather models can produce strong medium-range forecasts directly from reanalysis data \citep{lam2023graphcast}. On the probabilistic side, the strongest recent systems are often trained directly with CRPS-style objectives. FGN trains marginal probabilistic forecasts with CRPS and still recovers useful joint structure \citep{alet2025fgn}. AIFS-CRPS takes a similar proper-scoring-rule direction and shows very strong ensemble skill \citep{lang2026aifscrps}. More recently, U-Cast showed that a surprisingly simple U-Net backbone, trained first with MAE and then fine-tuned with CRPS using Monte Carlo Dropout, can be a very strong probabilistic weather forecaster \citep{cachay2026ucast}.

The issue is that optimizing forecast skill directly is not exactly the same thing as sampling physically clean weather trajectories. A model can score well in CRPS while still producing fields with odd spatial texture, poor balance, or unstable rollout behavior. Diffusion models attack the problem from the other direction: instead of predicting a distribution through a direct scoring rule, they learn to generate samples by iterative denoising. GenCast and cBottle are examples of this generative direction for weather and climate data \citep{price2025gencast,brenowitz2025cbottle}. These models tend to be good at sample quality and joint structure, but they are expensive because diffusion needs multiple denoising evaluations per forecast step \citep{ho2020ddpm,salimans2022progressive}.

The core idea of this project is to combine the two sides. I will use the U-Cast methodology as the probabilistic backbone to get an initial forecast, then run a diffusion refiner starting from that forecast. In other words: do not ask diffusion to generate the weather state from pure noise. Start it near a plausible forecast, add controlled noise, and let diffusion clean it up. This follows the basic idea of warm-start diffusion \citep{scholz2025warmstarts}, but applies it to low-resolution probabilistic weather forecasting.

\subsection{Related Work}

There are three main pieces of related work.

First, neural weather forecasting models give us the general forecasting setup. GraphCast showed that learned neural models can be competitive for medium-range global forecasting \citep{lam2023graphcast}. I are not trying to reproduce the full GraphCast system here, but it is still useful as background because the problem setting is similar: predict future global atmospheric states from recent reanalysis states.

Second, CRPS-trained probabilistic models give us the warm start. FGN is the closest conceptual reference point for direct probabilistic training: it generates probabilistic forecasts through learned model perturbations and trains directly on marginal CRPS \citep{alet2025fgn}. AIFS-CRPS is another strong example of using a CRPS-based objective to train an ensemble weather model efficiently \citep{lang2026aifscrps}. U-Cast is the most direct backbone reference for this project. It uses a standard U-Net architecture, deterministic pretraining, CRPS fine-tuning, and Monte Carlo Dropout to get a fast probabilistic forecaster \citep{cachay2026ucast}. I will use this U-Cast-style model as the first-stage probabilistic model, due to its simplicity and compatibility with Pytorch.

Third, diffusion models give us the sample refiner. Standard diffusion models learn a denoising process that maps noise to samples \citep{ho2020ddpm}. In weather, GenCast uses conditional diffusion for global probabilistic forecasting \citep{price2025gencast}, while cBottle uses a generative diffusion framework for high-resolution climate emulation and related tasks \citep{brenowitz2025cbottle}. Warm-start diffusion shows that conditional generation can be accelerated by replacing the uninformed noise prior with an informed, context-dependent starting point \citep{scholz2025warmstarts}. Our project is basically that idea, but with the warm start coming from a U-Cast-style probabilistic weather model.

\subsection{High-level Methodology}

I will build a two-stage probabilistic forecasting system.

The first stage is a U-Cast-style probabilistic forecaster. The input will be recent ERA5 states on the $5^\circ$ grid, and the output will be a probabilistic prediction for the next forecast step or a short rollout. The model will use a U-Net backbone. I will follow the U-Cast training recipe: first train the U-Net deterministically with a pointwise loss such as MAE, then fine-tune it probabilistically with a CRPS objective. Stochasticity will come from Monte Carlo Dropout or a closely related perturbation mechanism.

The second stage is a warm-start diffusion refiner. Importantly, the diffusion model will use the same basic U-Net architecture as the probabilistic backbone. Given the current weather context $c$ and a first-stage probabilistic forecast $\hat{x}_{t+\Delta}$, I create a noisy warm start
\[
z_\tau = \hat{x}_{t+\Delta} + \sigma_\tau \epsilon,
\qquad \epsilon \sim \mathcal{N}(0, I).
\]
A conditional U-Net denoiser then takes $z_\tau$, the weather context $c$, the backbone forecast $\hat{x}_{t+\Delta}$, and the noise level $\tau$ as input.

I will train two diffusion variants with different objectives. The first will use the standard diffusion noise-prediction objective (with latitude weighting for Climate), where the U-Net predicts the injected noise $\epsilon$. The second will use a more forecast-specific correction objective, where the U-Net predicts either the clean target $x_{t+\Delta}$ or the residual $x_{t+\Delta} - \hat{x}_{t+\Delta}$. The point is to test whether warm-start diffusion works better as a generic denoising model or as a residual physical correction model.

At inference time, we sample an initial forecast from the U-Cast-style probabilistic model, add noise at a chosen warm-start level, and run only a small number of denoising steps. The diffusion model is therefore not responsible for the entire forecast distribution. It only has to correct and sharpen an already reasonable forecast.

We can compare four systems:
\begin{enumerate}
    \item a deterministic U-Net trained with MAE,
    \item a U-Cast-style probabilistic U-Net trained with CRPS,
    \item a cold-start conditional diffusion U-Net,
    \item the proposed warm-start diffusion U-Net.
\end{enumerate}

Evaluation will focus on both forecast skill and sample quality. For probabilistic skill, we will report CRPS, ensemble spread-error consistency, and reliability-style diagnostics. For deterministic skill, I will report RMSE. For physical/sample quality, we can inspect spatial spectra and rollout stability.

\bibliographystyle{plainnat}
\bibliography{refs}

\end{document}